\section{中国剩余引理的一些推论}\label{sec:galois-1-9}
本节中，$A$ 是一般的交换环.
\subsection{引理}\label{sec:galois-1-9-1}
\begin{lemma*}
设理想 $I=\mathfrak m_1^{v_1}\mathfrak m_2^{v_2}\cdots\mathfrak m_r^{v_r}$，其中 $\mathfrak m_1,\ldots,\mathfrak m_r$ 是两两不同的极大理想. 则典范投影的乘积给出 $A$-代数同构
\[A/I\xrightarrow{\sim}A/\mathfrak m_1^{v_1}\times\cdots\times A/\mathfrak m_r^{v_r}.\]
\end{lemma*}
\begin{proof}
我们已经证明了如下形式的中国剩余定理：
\[J+K=A\quad\Longrightarrow\quad A/(J\cap K)\xrightarrow{\sim}A/J\times A/K.\]
注意，$JK$ 总是包含于 $J\cap K$；当 $J+K=A$ 时两者相等，因为此时
\[J\cap K=(J\cap K)J+(J\cap K)K\subset JK.\]
所以也可以表述为
\[J+K=A\quad\Longrightarrow\quad A/JK\xrightarrow{\sim}A/J\times A/K.\]
现证 $\mathfrak m_1^{v_1}+(\mathfrak m_2^{v_2}\cdots\mathfrak m_r^{v_r})=A$. 这将给出
\[A/I\simeq A/\mathfrak m_1^{v_1}\times A/(\mathfrak m_2^{v_2}\cdots\mathfrak m_r^{v_r}),\]
然后对 $r$ 归纳即可得出引理. 因为 $\mathfrak m_1$ 与 $\mathfrak m_i$（$i\ne1$）是不同的极大理想，对 $i>1$ 有 $\mathfrak m_1+\mathfrak m_i=A$. 从而
\begin{align*}
A&=(\mathfrak m_1+\mathfrak m_i)^{v_i}\\
&\subset\mathfrak m_1^{v_i}+\mathfrak m_1^{v_i-1}\mathfrak m_i+\cdots+\mathfrak m_1\mathfrak m_i^{v_i-1}+\mathfrak m_i^{v_i}
\subset\mathfrak m_1+\mathfrak m_i^{v_i}.
\end{align*}
因此 $A=\mathfrak m_1+\mathfrak m_i^{v_i}$. 同样，$A=(\mathfrak m_1+\mathfrak m_i^{v_i})^{v_1}\subset\mathfrak m_1^{v_1}+\mathfrak m_i^{v_i}=A$. 最后，
\begin{align*}
A&=(\mathfrak m_1^{v_1}+\mathfrak m_2^{v_2})(\mathfrak m_1^{v_1}+\mathfrak m_3^{v_3})\cdots(\mathfrak m_1^{v_1}+\mathfrak m_r^{v_r})\\
&\subset\mathfrak m_1^{v_1}+(\mathfrak m_2^{v_2}\mathfrak m_3^{v_3}\cdots\mathfrak m_r^{v_r}),
\end{align*}
这正是所需结论.
\end{proof}
\paragraph{记号.} 对 $A$-模 $M$ 和 $A$ 的理想 $I$，记
\begin{align*}
M[I]&:=\{m\in M:\ \forall i\in I,\ im=0\},\\
M[I^\infty]&:=\bigcup_{n\in\mathbb N}M[I^n]\\
&=\{m\in M:\ \exists n\in\mathbb N,\ \forall i_1,\ldots,i_n\in I,\ i_1\cdots i_nm=0\}.
\end{align*}
它们都是 $M$ 的 $A$-子模. 若 $IM=0$，也就是说 $M[I]=M$，就称 $M$``被 $I$ 零化''. 另回顾，$\operatorname{Max}(A)$ 表示 $A$ 的极大理想集合.
\subsection{定理}\label{sec:galois-1-9-2}
\begin{theorem*}
设 $A$ 是交换环，$M$ 是被有限个极大理想的乘积零化的 $A$-模.
\begin{enumerate}[label=\roman*)]
\item 对每个 $\mathfrak m\in\operatorname{Max}(A)$，有 $M[\mathfrak m^\infty]\simeq M_{\mathfrak m}$，且集合 $\{\mathfrak m\in\operatorname{Max}(A):M_{\mathfrak m}\ne0\}$ 有限.
\item 有 $M=\bigoplus_{\mathfrak m\in\operatorname{Max}(A)}M[\mathfrak m^\infty]\simeq\prod_{\mathfrak m}M_{\mathfrak m}$.
\end{enumerate}
\end{theorem*}
\begin{proof}
由假设，$M$ 被形如 $I=\mathfrak m_1^{v_1}\cdots\mathfrak m_r^{v_r}$ 的理想零化，其中 $\mathfrak m_i$ 是两两不同的极大理想. 由上面的引理及等式 $IM=0$，得到 $M$ 的分解
\begin{align*}
(*)\qquad M=M/IM&\simeq(A/I)\otimes_A M\\
&\simeq\prod_{i=1}^r(A/\mathfrak m_i^{v_i})\otimes_A M
=\prod_{i=1}^r M/\mathfrak m_i^{v_i}M.
\end{align*}
现注意，若 i)、ii) 对具有有限支撑的 $M$ 和 $M'$ 成立，它们就对 $M\oplus M'$ 成立，而后者显然也具有有限支撑. 根据分解 $(*)$，可以假设 $I=\mathfrak n^v$，其中 $\mathfrak n\in\operatorname{Max}(A)$.

先设 $\mathfrak m\ne\mathfrak n$. 前一引理中已证明 $\mathfrak m+\mathfrak n^v=A$，故可取 $p\in\mathfrak m$、$q\in\mathfrak n^v$ 使 $p+q=1$. 此时 $p\in\mathfrak m$ 在 $M$ 上的作用是恒等映射，故 $M[\mathfrak m^\infty]=0$. 此外，$q$ 零化 $M$，而 $q\in A\setminus\mathfrak m$，所以 $M_{\mathfrak m}=0$. 这样 $M[\mathfrak m^\infty]$ 与 $M_{\mathfrak m}$ 确实同构，而且都为零；这也说明 i) 中的集合至多有一个元素.

现设 $\mathfrak m=\mathfrak n$. 因为 $\mathfrak n^v$ 零化 $M$，有 $M=M[\mathfrak n^\infty]=M[\mathfrak n^v]=M/\mathfrak n^vM$. 为证明 $M\xrightarrow{\sim}M_{\mathfrak n}$，只需说明每个 $x\in A\setminus\mathfrak n$ 都在 $M$ 上可逆地作用. 由 $\mathfrak n$ 的极大性，$(x)+\mathfrak n=A$. 与前一引理的证明一样，得到 $(x)+\mathfrak n^v=A$，故存在 $y\in A$、$q\in\mathfrak n^v$ 使 $xy+q=1$. 因 $q$ 零化 $M$，$x$ 在 $M$ 上的作用可逆，其逆是 $y$ 的作用.
\end{proof}
\begin{remark*}
若 $M$ 被 $\mathfrak m_1^{v_1}\cdots\mathfrak m_r^{v_r}$ 零化，其中 $\mathfrak m_i$ 是两两不同的极大理想，则定理中的分解更具体地写成
\[M=M[\mathfrak m_1^{v_1}]\oplus\cdots\oplus M[\mathfrak m_r^{v_r}]\simeq M_{\mathfrak m_1}\times\cdots\times M_{\mathfrak m_r}.\]
\end{remark*}
\subsection{有限长度模}\label{sec:galois-1-9-3}
下面给出满足上述定理假设的一个重要例子.
\begin{definition*}
称 $A$-模 $M$ 为
\begin{itemize}
\item \newterm{galois-simple-module}{单模}{simple module}，若它非零，且仅有子模 $\{0\}$ 和 $M$；
\item \newterm{galois-finite-length-module}{有限长度模}{module of finite length}，若它有有限严格递增子模列 $0=M_0\subsetneq M_1\subsetneq\cdots\subsetneq M_n=M$，使每个 $M_i/M_{i-1}$（$i=1,\ldots,n$）都是单模.
\end{itemize}
\end{definition*}
若 $M$ 是 $A$-模，元素 $m\in M$ 的\newterm{galois-annihilator}{零化子}{annihilator}是子集
\[\operatorname{Ann}_A(m):=\{a\in A:am=0\}.\]
它是 $A$ 的理想. 模 $M$ 的零化子 $\operatorname{Ann}_A(M)$ 则是理想
\[\operatorname{Ann}_A(M):=\bigcap_{m\in M}\operatorname{Ann}_A(m)=\{a\in A:aM=0\}.\]
\begin{lemma*}
单模的零化子是极大理想. 有限长度模的零化子是有限个极大理想的乘积.\footnote{译注：原文如此表述；紧随的证明只说明存在有限个极大理想的乘积零化该模，即其零化子包含这样的乘积.}
\end{lemma*}
\begin{proof}
若 $M$ 是单模，$m\in M\setminus\{0\}$，则 $Am=M$，且同态 $A\longrightarrow M$、$a\mapsto am$ 诱导同构 $A/\operatorname{Ann}_A(m)\xrightarrow{\sim}M$. 因 $M$ 是单模，商环 $A/\operatorname{Ann}_A(m)$ 没有非零真理想，故为域，因而 $\operatorname{Ann}_A(m)$ 是极大理想. 现设 $M$ 为任意 $A$-模. 注意，对于 $M$ 的子模 $N$，若 $I$ 零化 $N$ 且 $J$ 零化 $M/N$，则 $IJ$ 零化 $M$. 事实上，$j\in J$ 在 $M/N$ 上的作用为零，故 $jM\subset N$，从而 $ijM\subset iN=0$. 直接归纳可知，若有滤过 $0=M_0\subsetneq M_1\subsetneq\cdots\subsetneq M_n=M$，且 $I_i$ 零化 $M_i/M_{i-1}$（$i=1,\ldots,n$），则 $I_1I_2\cdots I_n$ 零化 $M$. 当 $M$ 有限长度时，由此得到陈述中的第二个断言.
\end{proof}
\begin{remark*}
映射 $M\mapsto\operatorname{Ann}_A(M)$ 与 $\mathfrak m\mapsto A/\mathfrak m$ 在单模的同构类与极大理想之间建立双射.
\end{remark*}
\begin{corollary*}
第~\ref{sec:galois-1-9-2}~分目的定理适用于每个有限长度模.
\end{corollary*}
\begin{example*}
每个有限 Abel 群 $M$ 都是有限长度的 $\mathbb Z$-模（例如可对基数归纳证明）. 因此有限 Abel 群典范地是其极大 $p$-子群的有限直积 $M=\prod_p M[p^\infty]$.
\end{example*}
\subsection{在有限维代数上的应用}\label{sec:galois-1-9-4}
下面是上述推论的一个重要特例. 设 $A$ 是域 $K$ 上的代数，$M$ 是有限维的 $A$-模. 则 $M$ 有限长度. 事实上，每个严格递增列 $0=M_0\subsetneq M_1\subsetneq\cdots\subsetneq M_n=M$ 的``长度'' $n$ 都不超过 $\dim_K(M)$，故可以选取``长度''最大的一个；这样的列的相继商必然是单模.
\begin{corollary*}
设 $A$ 是域 $K$ 上的有限维交换代数. 则 $\operatorname{Max}(A)$ 有限，且 $A$ 是其在各个极大理想处的局部化的直积：
\[A\xrightarrow{\sim}\prod_{\mathfrak m\in\operatorname{Max}(A)}A_{\mathfrak m}.\]
\end{corollary*}
\begin{proof}
这是上述结论中 $M=A$ 的情形.
\end{proof}
\begin{example*}
设 $A=\mathbb C[X]/(f)$，其中 $f$ 首一. 将其分解为 $f=\prod_{i=1}^r(X-z_i)^{v_i}$. 中国剩余定理给出
\[A\xrightarrow{\sim}\mathbb C[X]/(X-f_1)^{v_1}\times\cdots\times\mathbb C[X]/(X-f_r)^{v_r}.\]
$A$ 的极大理想是 $\mathfrak m_i$（$i=1,\ldots,r$），分别由 $X-z_i$ 在 $A$ 中的像生成，而局部化 $A_{\mathfrak m_i}$ 正是因子 $\mathbb C[X]/(X-f_i)^{v_i}$.\footnote{译注：原文的分解因子及末式使用 $f_i$，而此前定义的根为 $z_i$；此处保留原记号.}
\end{example*}
回顾，$A$-代数 $A_{\mathfrak m}$ 只有一个极大理想，即 $\mathfrak m A_{\mathfrak m}$；称它是局部的. 上述推论因而说明，域上的每个有限维代数都是局部代数的直积. 对这种代数的结构还可以作进一步的说明.
\begin{proposition*}
设 $A$ 是域 $K$ 上的有限维局部代数，$\mathfrak m$ 是它的极大理想. 则 $\mathfrak m$ 幂零. 更确切地说，若 $d=\dim_K(A)$，则 $\mathfrak m^d=0$.
\end{proposition*}
\begin{proof}
递减列 $\mathfrak m\supset\mathfrak m^2\supset\cdots\supset\mathfrak m^n\supset\cdots$ 在指标 $n=d$ 以前就稳定，因为这些都是有限维 $K$-向量空间. 设 $\mathfrak m^{r+1}A=\mathfrak m^rA$. 下面著名的引理保证 $\mathfrak m^rA=0$.
\end{proof}
\begin{lemma*}[{Nakayama 引理}]
设 $A$ 是极大理想为 $\mathfrak m$ 的局部环，$M$ 是有限生成 $A$-模. 若 $M=\mathfrak m M$，则 $M=0$.
\end{lemma*}
\begin{proof}
设 $M\ne0$，取 $M$ 的一个极小生成集 $\{m_1,\ldots,m_r\}$. 因为 $m_r\in\mathfrak m M$，存在 $i_1,\ldots,i_r\in\mathfrak m$ 使 $m_r=i_1m_1+\cdots+i_rm_r$. 但 $1-i_r\notin\mathfrak m$，故理想 $(1-i_r)$ 不包含于 $\mathfrak m$；由于 $\mathfrak m$ 是唯一极大理想，它不是真理想，因而等于 $A$，所以 $1-i_r\in A^\times$. 从而
\[m_r=(1-i_r)^{-1}(i_1m_1+\cdots+i_{r-1}m_{r-1}),\]
这与所取生成集的极小性矛盾.
\end{proof}
\paragraph{应用.} 域 $K$ 上的有限维既约交换代数 $A$ 是 $K$ 的有限扩张的有限直积 $A=K_1\times\cdots\times K_r$.

